ECFP is Weisfeiler--Leman colour refinement followed by hashing, so folding a fingerprint at
radius $r$ into $B$ bits incurs exactly two losses, and both can be measured directly. The
\emph{truncation} loss is the fraction of WL-1 colour classes still unresolved after $r$ rounds,
relative to the stable partition obtained by refining to convergence; we compute it by running
refinement on the molecular graph, using the same initial invariant as the fingerprint
(atomic number, degree, attached hydrogens, formal charge, isotope, ring membership, chiral tag),
until the partition stops becoming finer. The \emph{collision} loss is
$1 - (\text{on-bits})/(\text{distinct environments})$, the fraction of distinct environments that
share a bit after folding. Table~\ref{tab:radius} reports both, pooled over 20{,}000 molecules
drawn from our corpus (22.2 heavy atoms mean), together with the fraction of molecules whose
WL-1 partition is already stable at that radius.

\begin{table}[h]
\centering
\begin{tabular}{crrrr}
\toprule
& \multicolumn{2}{c}{truncation} & \multicolumn{2}{c}{collision} \\
\cmidrule(lr){2-3}\cmidrule(lr){4-5}
$r$ & loss & WL-stable & 2048 bits & 4096 bits \\
\midrule
1 & 22.50\% & 12.0\% & 1.21\% & 0.14\% \\
2 &  7.75\% & 52.9\% & 1.34\% & 0.35\% \\
\textbf{3} & \textbf{3.23\%} & \textbf{80.0\%} & \textbf{1.61\%} & 0.56\% \\
4 &  1.63\% & 90.6\% & 1.85\% & 0.71\% \\
5 &  0.87\% & 95.6\% & 2.02\% & 0.82\% \\
\bottomrule
\end{tabular}
\caption{The two losses of a folded Morgan fingerprint relative to exact WL-1 refinement, pooled
over 20{,}000 molecules. Truncation falls steeply with radius while collisions rise slowly, and
$r = 3$ is where the two become comparable.}
\label{tab:radius}
\end{table}

The two losses run in opposite directions, which is what makes the radius a genuine trade-off
rather than a convention. Truncation falls steeply --- a radius-1 fingerprint leaves more than a
fifth of the WL-1 partition unresolved and only 12\% of molecules fully refined --- while
collisions rise slowly, since a wider radius adds environments faster than it exhausts the bit
space. Moving from $r = 2$ to $r = 3$ removes 4.5 points of truncation at a cost of 0.27 points
of collision, a return of roughly 17 to 1, and takes the fraction of molecules whose colour
partition is already stable from 53\% to 80\%. The next step, $r = 3$ to $r = 4$, returns only
6.7 to 1, and $r = 4$ to $r = 5$ just 4.5 to 1. We therefore adopt $r = 3$: it is the last radius
at which truncation is the dominant loss, and beyond it the fingerprint pays increasing collision
for distinctions most molecules no longer have. Widening to 4096 bits would cut the collision
term by roughly two thirds at every radius, but doubles the memory of the representation for a
loss that is already the smaller of the two, so we retain 2048 bits.

Note that both quantities are representational rather than predictive. They bound how much of the
WL-1 colouring the fingerprint fails to carry; they do not establish that the missing
distinctions matter for any particular endpoint.
